% Inclusion-ready fragment. Required packages: amsmath, amssymb, amsthm, hyperref.
% Bibliography keys are supplied by nonstandard_rh_sources.bib.
\section{The nonstandard Riemann-hypothesis bridge in the original GCT sources}
\label{sec:gct-nrh-source-bridge}

The bridge named by the user is explicitly present in Mulmuley's
\emph{On P vs. NP, Geometric Complexity Theory, and the Riemann Hypothesis},
arXiv:0908.1936v2~\cite{mulmuley2009gctrh}.  Its final section, headed ``How to prove PH1 and why
should it hold?'', proposes extending the geometric ingredients surrounding
the Riemann hypothesis over finite fields to nonstandard quantum groups.
The stated destination is a characteristic-zero complexity separation.
This is an actual research programme in the source.  In that source the
nonstandard Riemann hypotheses themselves are described as still unknown;
no nonstandard zeta function, Frobenius space, weight bound, or numbered
nonstandard-RH conjecture is specified there.  Consequently the reconstruction
below retains the source's proved constructions and its conjectured arrows
with their original status.  It does not replace the proposed finite-field
bridge by the classical hypothesis for the complex zeros of
$\zeta_{\mathbb Q}(s)$.

\subsection{Objects and arrows before any complexity claim}

In GCT VII, arXiv:0709.0749v2~\cite{mulmuley2008gct7}, Section~1, $H$ is a complex connected
classical reductive group, $\mu$ is a highest weight,
$X=V_\mu(H)$, $G=GL(X)$, and
\[
 \rho:H\longrightarrow G
\]
is the representation homomorphism.  The plethysm multiplicity is
\[
 a_{\lambda,\mu}^{\pi}
 =\dim_{\mathbb C}\operatorname{Hom}_{H}
       \bigl(V_\pi(H),\operatorname{Res}_{\rho}V_\lambda(G)\bigr).
\]
Here $\lambda$ is a highest weight of $G$ and $\pi$ a highest weight of
$H$.  This notation is that of GCT VII; the survey uses a different order
of indices in its plethysm discussion, which is not silently substituted
in these formulas.

The quantum parameter $q$ is independent of any prime-power cardinality.
GCT VII's Section~6, \texttt{srefined}, explicitly warns that a splitting
extension $K/\mathbb Q(q)$ may be needed.  Its general reciprocity statement
is not asserted over $\mathbb Q(q)$ without that extension, and Section~7.1
gives a failure over $\mathbb Q(q)$.  The following description therefore
retains the chosen coefficient field $K$, rather than erasing it.

Let $X_q=V_{q,\mu}(H_q)$ and let $\widehat R_{X,X}^{H}$ be the braiding
operator attached to the standard quantum group $H_q$.  GCT VII,
Section~2, splits $X_q\otimes_K X_q$ according to the signs of the
eigenvalues $\pm q^{a/2}$ and writes
\[
 I=P_{X,X}^{+,H}+P_{X,X}^{-,H}.
\]
With $u^i_j$ the entries of a variable $\dim(X)\times\dim(X)$ matrix,
its matrix bialgebra is the explicit quotient
\begin{equation}\label{eq:gct-nrh-frt-source}
 \mathcal O(M_q^H(X))
 =K\langle u^i_j\rangle\Big/
 \left\langle
 P_{X,X}^{+,H}(\mathbf u\otimes\mathbf u)
       -(\mathbf u\otimes\mathbf u)P_{X,X}^{+,H}
 \right\rangle.
\end{equation}
Replacing $P^{+,H}$ by $P^{-,H}=I-P^{+,H}$ gives the same ideal because
the corresponding matrix of relations changes by the scalar $-1$.
The source defines braided symmetric and exterior algebras with the
relations $P^{-,H}\mathbf x_1\mathbf x_2=0$ and
$P^{+,H}\mathbf x_1\mathbf x_2=0$, respectively.  Nonstandard minors
are the matrix coefficients of the resulting exterior-algebra coactions.
The top exterior power can vanish.  Thus the determinant localization
and an antipode cannot be presumed to exist in the general plethysm case.
The source's term ``quantum group'' includes this possible semigroup case.

Theorems labelled \texttt{tnonstd} and \texttt{tnonstdformal} in GCT VII
state the construction of the quantum map $H_q\to G_q^H$, its matrix
coordinate bialgebra map, cosemisimplicity, and a Peter--Weyl decomposition.
The coordinate-ring direction is
\[
 \mathcal O(M_q^H(X))\longrightarrow\mathcal O(M_q(V))
 \qquad(H=GL(V)),
\]
which is opposite to the displayed direction $H_q\to G_q^H$ of virtual
quantum spaces.  We use the concrete coordinate map, whose elementary
quotient construction is proved below, without importing an unproved
canonical basis from the existence of the bialgebra.

GCT VII's conjecture \texttt{cduality} proposes a decomposition
\begin{equation}\label{eq:gct-nrh-duality-source}
 X_q^{\otimes r}
   \cong\bigoplus_\alpha W_{q,\alpha}\otimes_K T_{q,\alpha}
\end{equation}
for the commuting nonstandard quantum-group and
$\mathcal B_r^H(q)$ actions.  Its conjecture \texttt{creciprocity}, refined
by \texttt{crecirefined}, proposes
\begin{equation}\label{eq:gct-nrh-reciprocity-source}
 V_{q,\lambda}^{H}=\bigoplus_\alpha
                 (W_{q,\alpha})^{\oplus m_\lambda^\alpha},
 \qquad
 V_{q,\lambda}^{H}\big|_{q=1}
       \cong\operatorname{Res}_{\rho}V_\lambda(G).
\end{equation}
The multiplicity $m_\lambda^\alpha$ counts the factors isomorphic to the
Specht module $S_\lambda$ in a composition series of
$T_{q,\alpha}(1)$ as a $\mathcal B_r^H(1)$-module.  A semisimple generic
algebra does not license replacing that composition series by an assumed
semisimple special fibre.  The source itself makes this point in
Section~6.  If $n_\pi^\alpha$ denotes the multiplicity of $V_{q,\pi}$ in
$W_{q,\alpha}$, the proposed reciprocity yields the source's equation
\begin{equation}\label{eq:gct-nrh-multiplicities-source}
 a_{\lambda,\mu}^{\pi}
       =\sum_\alpha m_\lambda^\alpha n_\pi^\alpha.
\end{equation}
This displayed equality is reported here as the consequence asserted by
that reciprocity conjecture, not as a newly proved general-plethysm result.
The elementary dimension calculation involved is explicit: on a direct
sum as in \eqref{eq:gct-nrh-reciprocity-source}, an $H_q$-linear map from
$V_{q,\pi}$ is exactly one such map to each summand, and projection and
inclusion supply inverse linear maps between the corresponding Hom spaces.
Each of the $m_\lambda^\alpha$ copies contributes
$n_\pi^\alpha$ dimensions.  What this calculation does not establish is
the asserted specialization isomorphism in
\eqref{eq:gct-nrh-reciprocity-source}.

\subsection{The positivity and effectivity steps retained separately}

GCT VIII, arXiv:0709.0751v2~\cite{mulmuley2008gct8}, presents algorithms whose correctness and
required properties are conjectured.  Its source labels include
\texttt{clocalcrystalb} (local crystal-basis existence),
\texttt{cuniquemindegleft} and \texttt{cunimindegright} (minimum degree),
\texttt{ccelldecompleft} and \texttt{ccelldecompright} (cell decomposition),
and \texttt{cposkroneckerleft}, \texttt{cposkroneckerright} (positivity).
The cells must recover the irreducible modules, and the subcells must
recover the restricted $H_q$-modules.  This is structural data beyond a
list of positive numbers.

In the Kronecker setting
\[
 H=GL(V)\times GL(W),\qquad X=V\otimes W,
\]
the left positivity conjecture in Section~2.4.2 does not assert that
every raw multiplication constant is a nonnegative polynomial.  It allows
the exact form
\begin{equation}\label{eq:gct-nrh-signed-positivity}
 \pm(q-q^{-1})^a f(q),\qquad a\in\mathbb Z_{\ge0},
\end{equation}
where $f$ is invariant under $q\leftrightarrow q^{-1}$ and has
nonnegative rational coefficients, together with the stated unimodality
requirement.  The factor $(q-q^{-1})^a$ and its sign are part of the
conjectured statement and cannot be suppressed at $q=1$.

The general-plethysm proposal is weaker and more complicated.  The
coefficients can be algebraic over $\mathbb Q(q)$, and Section~2.4.3
discusses coefficients of their minimal polynomials.  The conjecture
\texttt{cnonstdposleft0} uses positivity of a nonzero residual polynomial
at $q=1$ after explicitly retaining signs and powers of $q-q^{-1}$.
The conjecture \texttt{cnonstdposleft1} proposes a decomposition
$c(q)=c^0(q)+c^1(q)$ with a dominant positive unimodal term and a small
correction.  The source includes asymptotic language and a polynomial
bitlength bound; it does not provide a fully specified Frobenius complex
whose weights prove that decomposition.

Section~6 of GCT VIII then adds computational requirements on the labels
of the basis and its cells.  In the standard setting, efficiently
computable tableau/path labels and crystal operations yield counting
rules even when constructing all basis vectors is expensive.  The
nonstandard analogues are proposed there as additional properties whose
full specification is deferred.  Thus coefficient positivity alone
does not provide a uniform polynomial-time algorithm or a $\#P$ formula.
For example, from a predicate $R(x,y)$ and a bound $|y|\le p(|x|)$, both
explicitly supplied, one obtains a $\#P$ count by nondeterministically
choosing a padded witness and checking $R$.  Without a witness encoding,
polynomial bound, and polynomial-time checking procedure, the mere
nonnegative integer value of a multiplicity supplies no such machine.

The survey's final section next discusses Plethysm PH0, the existence
of compatible positive bases for $\rho:H\to GL(V_\mu(H))$, and
Plethysm PH/PH1, positive polyhedral formulas and their effective
versions.  The intended further step is to quantize the triple
\[
 \overline H_{\widehat h}\longrightarrow
 \overline H=SL(\overline X)\longrightarrow L=GL(\overline W)
\]
associated with a class variety.  It is not a proved automatic passage
from a single plethysm pair to all class-variety multiplicities.
The survey's general PH1, labelled \texttt{hph1restatement}, requires
bounds polynomial in $\langle\lambda\rangle,n,\langle m\rangle$,
including the bitlength of $m$, for the padded polynomial
$\overline h(\overline X)=z^{m-n}h(X)$.

Finally the displayed source theorem \texttt{tgct6refined} also requires
an obstruction hypothesis OH.  It takes
$m=2^{\log^a n}$, with $a>1$ fixed, and asks for labels $\lambda$ with
$P_{\lambda,n,m}(k)\ne\varnothing$ for all sufficiently large $k$
but $Q_\lambda=\varnothing$.  This separates a positive formula from
the additional proof that the required obstruction actually exists.
Its displayed target is the general $\#P$ versus $NC$ problem in
characteristic zero; the source routes the analogous $P$ versus $NP$
formulation to GCT VI.  This record does not substitute a Boolean uniform
$P$ versus $NP$ theorem for that stated algebraic target.

The primary-source chain therefore has the following exact status:
\begin{enumerate}
\item The standard geometric positivity story uses the finite-field
purity results surrounding Deligne's work and the canonical bases of
standard quantum groups.  This is the source's established model.
\item The matrix bialgebra quantization is explicitly constructed.
\item General nonstandard canonical bases, their required cells,
positivity, reciprocity and effective labels are proposed additional
statements.  They are not all consequences of the existence of the
matrix bialgebra.
\item The nonstandard finite-field purity/RH theory is a proposed way to
prove those statements; the cited works do not state it as a complete
zeta-function hypothesis.
\item Effective positive formulas must still be connected to the class
varieties and to actual obstruction existence, with the declared size
parameters, before the intended complexity separation follows.
\end{enumerate}
These labels describe the literature's proposed programme, and introduce
no conditional theorem as a replacement for a missing proof in the
present work.

\subsection{A fully explicit FRT coordinate morphism}

Here is the elementary algebraic construction that underlies
\eqref{eq:gct-nrh-frt-source}.  It is included with a full proof so that
the quantum-space arrow has a concrete coordinate-ring meaning.
Let $K$ be a field, $E=K^d$ with its ordered basis $e_1,\ldots,e_d$, and
$P\in\operatorname{End}_K(E\otimes E)$.  Put
\[
 T=K\langle u_{ij}:1\le i,j\le d\rangle,
 \quad W_{(i,k),(j,l)}=u_{ij}u_{kl},
 \quad R=PW-WP,
 \quad A_P=T/(R_{ab}:a,b\in\{1,\ldots,d\}^2).
\]
Here $(R_{ab})$ denotes the two-sided ideal generated by the displayed
entries; no commutation between different $u_{ij}$ is presumed.
On the free algebra define
\[
 \Delta(u_{ij})=\sum_{h=1}^{d}u_{ih}\otimes u_{hj},
 \qquad \epsilon(u_{ij})=\delta_{ij}.
\]
Both extend uniquely as algebra homomorphisms.  Their coassociativity
and counit identities hold on a generator because both iterated
coproducts equal
$\sum_{h,k}u_{ih}\otimes u_{hk}\otimes u_{kj}$, and contracting either
factor with $\epsilon$ gives $u_{ij}$.  The same identities then hold on
each word by multiplicativity and on all of $T$ by linearity.

For indices $a,b,c$ in $\{1,\ldots,d\}^2$,
\[
 \Delta(W_{ab})=\sum_c W_{ac}\otimes W_{cb}.
\]
Consequently, with each occurrence of the scalar matrix $P$ acting on
these pair-indices,
\[
 \Delta(R_{ab})
 =\sum_c R_{ac}\otimes W_{cb}
       +\sum_c W_{ac}\otimes R_{cb},
 \qquad \epsilon(R_{ab})=0.
\]
The two middle terms cancel entry by entry.  The ideal of relations is
therefore a coideal and lies in the kernel of the counit.  It follows
that these formulas define a bialgebra structure on $A_P$.

Now let $B$ be a bialgebra with a right comodule $E$ written
\[
 \delta_B(e_j)=\sum_i e_i\otimes y_{ij}.
\]
The comodule identities are exactly
$\Delta_B(y_{ij})=\sum_h y_{ih}\otimes y_{hj}$ and
$\epsilon_B(y_{ij})=\delta_{ij}$.  When $P$ is a comodule endomorphism
of $E\otimes E$, applying $\delta_B P=(P\otimes1)\delta_B$ to every
basis vector gives
$P(y\otimes y)=(y\otimes y)P$ entry by entry.  Thus the explicit map
\begin{equation}\label{eq:gct-nrh-universal-map}
 \Phi_P:A_P\longrightarrow B,\qquad
       [u_{ij}]\longmapsto y_{ij}
\end{equation}
is well-defined.  It preserves coproduct and counit on generators and
hence everywhere.  It is the unique bialgebra map with those specified
generator images.  Its image is the subbialgebra generated by the
$y_{ij}$; its exact kernel is
\[
 \{[F(u)]\in A_P:F(y)=0\text{ in }B\}.
\]
This does not claim that the map is injective, nor that its image is all
of $B$.

For comparison at the classical symmetric projector, assume
$\operatorname{char}(K)\ne2$ and take $P=(I+\tau)/2$, where
$\tau(e_i\otimes e_k)=e_k\otimes e_i$.  Then
\[
 2R_{(i,k),(j,l)}
   =u_{kj}u_{il}-u_{il}u_{kj}.
\]
As the four indices vary these are exactly all pairwise commutators of
the matrix variables.  Sending $[u_{ij}]$ to the corresponding commuting
variable defines an isomorphism
\[
 A_{(I+\tau)/2}\longrightarrow K[u_{ij}:1\le i,j\le d],
\]
whose inverse sends a commuting variable to $[u_{ij}]$.  The commutator
relations prove that the inverse is defined, and the two compositions
fix every generator.  This proves the comparison map with its field,
size and characteristic restriction retained.

\subsection{Specialization and the information in its kernel}

Let $D=\mathbb Q[q,q^{-1}]_{(q-1)}$, $t=q-1$, $L$ a finite free
$D$-module, $J\subseteq L$ a submodule, and $A=L/J$.  These are explicit
algebraic data rather than a claim that a particular GCT conjecture
holds.  Let $A[t]=\{a\in A:ta=0\}$.  There is an exact sequence
\begin{equation}\label{eq:gct-nrh-specialization-exact}
 0\longrightarrow A[t]\xrightarrow{\eta}J/tJ
 \longrightarrow L/tL\longrightarrow A/tA\longrightarrow0,
 \qquad \eta([v])=[tv].
\end{equation}
Indeed $ta=0$ means that a lift $v\in L$ satisfies $tv\in J$.
Changing the lift by $j\in J$ changes $tv$ by $tj$, so $\eta$ is
well-defined.  If $tv\in tJ$, then $tv=tj$ for some $j\in J$; multiplication
by $t$ is injective on the free module $L$, so $v=j$ and $[v]=0$.
Thus $\eta$ is injective.  An element $j\in J$ maps to zero in $L/tL$
exactly when $j=tv$ for some $v\in L$, which is exactly the image of
$\eta$.  The kernel of $L/tL\to A/tA$ consists of the classes of
$J+tL$, and the last map is surjective by lifting a representative from
$A$ to $L$.  These observations prove every assertion of exactness.

Put
\[
 J^{\mathrm{sat}}=
  \{v\in L:t^e v\in J\text{ for some }e\in\mathbb Z_{\ge0}\}.
\]
The map $A\to A\otimes_D\mathbb Q(q)$ has kernel
$J^{\mathrm{sat}}/J$.  To prove this, localization sends the class of
$v$ to zero precisely when some nonzero element of $D$ multiplies $v$
into $J$.  Every nonzero element of $D$ is a power of $t$ times a unit:
factor its numerator by its exact order of vanishing at $q=1$ and note
that its denominator and remaining numerator are units in this local
ring.  Multiplying or dividing by that unit does not change membership
in $J$.  This proves the kernel formula.  Passing to a generic fibre
therefore removes an explicitly identified module, not an unspecified
``small error''.  Formula \eqref{eq:gct-nrh-specialization-exact} records
the first special-fibre layer that a proposed reciprocity construction
must preserve in the data to which it is applied.

\subsection{A source-specific coefficient and its retained double zero}

The final GCT IV, by Blasiak, Mulmuley and Sohoni,
arXiv:cs/0703110v4~\cite{blasiakmulmuleysohoni2013gct4}, proves nonstandard Schur--Weyl duality in Section~12
(source label \texttt{t nonstandard schur-weyl duality}).  Its displayed
statement uses $\check X^{\otimes r}$, and the immediately following
remark explicitly records the dual convention needed by its proof.
The source then constructs a global crystal-basis rule for the
two-row problem, $\dim V=\dim W=2$, with
\[
 g_{\lambda\mu\nu}
  =\#\{\text{highest-weight elements of }{+}\operatorname{HNSTC}(\nu)
          \text{ of weight }(\lambda,\mu)\}.
\]
This is a proved bounded case in the source, not a claim about all
dimensions.  Theorem \texttt{t positivity} gives sign-coherent
coefficients for individual Chevalley generators.  Example
\texttt{ex not positive} displays a divided-power expression with mixed
signs. The literal printed target tableau has a weight inconsistency.
Section~\ref{sec:gct-generator-interval} proves that obstruction, constructs
the corrected target $Q$ from the actual action rules, and proves that
its divided-power matrix entry is the displayed Laurent polynomial.
The printed and corrected targets remain separately specified objects.

We now independently prove all algebraic identities used from that
example.  The example takes column-count shape
$\nu=[0,n_3,n_2,2]$ with $n_3=3,n_2=2$; in ordinary row-length notation
this is the same partition $(7,5,3)$, since there are three columns of
height three, two of height two and two of height one.  Its size is
$3\cdot3+2\cdot2+2=15$.  We retain the original column-count shape.
For $n\ge1$ set
\[
 [n]_q=\sum_{j=0}^{n-1}q^{n-1-2j}
 \quad\text{in }\mathbb Z[q,q^{-1}].
\]
Multiplication gives
$(q-q^{-1})[n]_q=q^n-q^{-n}$ by cancellation of explicitly adjacent
terms; the finite sum also defines its values at $q=1$ and $q=-1$.
The source displays the following rational expression and Laurent polynomial:
\begin{equation}\label{eq:gct-nrh-coefficient-original}
 C(q)=\frac{[6]_q\bigl([7]_q-[3]_q\bigr)-[3]_q[8]_q}{[2]_q}
      =q^{10}-q^4-q^{-4}+q^{-10}.
\end{equation}
The first expression lives in $\mathbb Q(q)$ and the second is its
Laurent-polynomial representative; multiplying the equality by
$[2]_q=q+q^{-1}$ verifies it without imposing any evaluation at a zero
of $[2]_q$.

The original Laurent polynomial, including all four signs, satisfies
\begin{align}
 C(q)
  &=(q^7-q^{-7})(q^3-q^{-3})\notag\\
  &=(q-q^{-1})^2[7]_q[3]_q.\label{eq:gct-nrh-coefficient-factors}
\end{align}
The first equality follows by multiplying its two binomials, giving
the four exponents $10,4,-4,-10$ with coefficients $1,-1,-1,1$.
The second follows from the preceding finite-sum identity for $n=7,3$.
Moreover
\[
 [7]_q[3]_q
 =q^8+2q^6+3q^4+3q^2+3+3q^{-2}+3q^{-4}+2q^{-6}+q^{-8}.
\]
The displayed coefficients count the ordered pairs of summands with
each exponent.  This proves nonnegativity of this retained factor,
while $C$ itself has the mixed signs displayed in
\eqref{eq:gct-nrh-coefficient-original}.  No assertion about every
positivity requirement of GCT VIII follows from this single example.

Since
\[
 q-q^{-1}=\frac{(q-1)(q+1)}q,
\]
equation \eqref{eq:gct-nrh-coefficient-factors} gives the equality in
$D$ with all factors retained:
\[
 C(q)=(q-1)^2\frac{(q+1)^2}{q^2}[7]_q[3]_q,
 \qquad
 \left.\frac{(q+1)^2}{q^2}[7]_q[3]_q\right|_{q=1}=4\cdot7\cdot3=84.
\]
Thus $C(1)=C'(1)=0$, $C''(1)=168$, and its exact order at $q=1$ is two.
Specializing $C$ to its value $0$ would lose the nonzero retained second
jet $84(q-1)^2$ in $D/(q-1)^3$.

There is also an explicit two-sheet presentation.  Here $p$ is a new trace
coordinate, distinct from the source polynomial map's retained coordinate
$w$. Define
\[
 B=\mathbb Q[p,Q]/(Q^2-pQ+1),\qquad
 \theta:B\longrightarrow\mathbb Q[q,q^{-1}],
 \quad Q\longmapsto q,\quad p\longmapsto q+q^{-1}.
\]
In $B$, $Q(p-Q)=1$.  Hence the inverse homomorphism sends $q$ to $Q$
and $q^{-1}$ to $p-Q$.  Both compositions fix the displayed generators,
proving that $\theta$ is an isomorphism.  The ring $B$ is free of rank
two over $\mathbb Q[p]$, with ordered basis $(1,Q)$: division by the
monic relation produces a unique representative $a(p)+b(p)Q$, and a
nonzero multiple of a monic degree-two polynomial cannot have smaller
$Q$-degree.  The deck involution sends $Q$ to $p-Q$; its image under
$\theta$ is $q\mapsto q^{-1}$.

Put $r_{\mathrm{GCT}}=2Q-p$ in $B$.  Then
\[
 r_{\mathrm{GCT}}^2=p^2-4.
\]
The associated conic presentation is the isomorphism
\[
 \mathbb Q[p,r_{\mathrm{GCT}}]/(p^2-r_{\mathrm{GCT}}^2-4)
  \longrightarrow\mathbb Q[q,q^{-1}],
 \quad (p,r_{\mathrm{GCT}})\longmapsto
        (q+q^{-1},q-q^{-1}),
\]
whose inverse sends $q$ to $(p+r_{\mathrm{GCT}})/2$ and $q^{-1}$ to
$(p-r_{\mathrm{GCT}})/2$.  Their product is $1$ by the conic equation.
The use of $1/2$ is explicit; these displayed inverses are over
$\mathbb Q$ and are not asserted in characteristic two.

Writing $p=q+q^{-1}$ and applying the recurrence
$[n+1]_q=p[n]_q-[n-1]_q$ gives
\[
 [3]_q=p^2-1,\qquad [7]_q=p^6-5p^4+6p^2-1.
\]
The recurrence is verified directly by multiplying the two finite
sums and cancelling the shared interior terms.  Thus the same retained
coefficient is
\begin{align*}
 C(q)
  &=(p^2-4)(p^6-5p^4+6p^2-1)(p^2-1)\\
  &=p^{10}-10p^8+35p^6-51p^4+29p^2-4.
\end{align*}
At $p=2$ the factor multiplying $p-2$ has value
$4\cdot7\cdot3=84$, whereas
\[
 p-2=\frac{(q-1)^2}{q}.
\]
This exact ramification equation explains why a first zero in the trace
coordinate $p$ is the retained second zero in the coordinate $q$.
The discriminant of $Q^2-pQ+1$ is $p^2-4$; the two excluded branch
values for the etale double cover are $p=2$ and $p=-2$.
Indeed the relative differential module is
$\Omega_{B/\mathbb Q[p]}=B/(2Q-p)\,dQ$, by differentiating the single
monic relation with $p$ held fixed.  Since $(2Q-p)^2=p^2-4$, this module
vanishes after inverting $p^2-4$.  The algebra is already free of rank
two over the base, so this localization is finite etale.  At $p=2$ the
fibre is $\mathbb Q[Q]/(Q-1)^2$, and at $p=-2$ it is
$\mathbb Q[Q]/(Q+1)^2$, proving that both exceptional fibres retain a
nonzero square-zero element.  Further, if $a(p)+b(p)Q$ is invariant under
$Q\mapsto p-Q$, then $b(p)(2Q-p)=0$.  The proved embedding into the
Laurent-polynomial domain gives $b=0$, so the invariant subring is
exactly $\mathbb Q[p]$.
These identities give a concrete, fully calculated nonstandard-GCT
coefficient and its rank-two coordinate algebra for any subsequent
arithmetic transport to act on.  They prove no general canonical-basis
conjecture and no complexity separation.
